\FloatBarrier
\subsection{Smooth flange revision}
\label{sec:ft-shield-smooth}

The shield/cone interference recorded below is resolved by the subsequent cone revision in
Section~\ref{sec:moller-cone}.

The smooth-flange FTInnerShield revision fills eight conical recesses, sixteen face chamfers, and two
stepped counterbores retained in the previous 13,500-triangle version. Both flange faces are complete flat
annuli. The large central passage, tube chamfers and flange fillets, and both longitudinal tube grooves
remain. The name, tungsten material, color, and all nine passive placements remain unchanged.

Removing the recess cuts reduces the reconstructed reference to 8,536 triangles. Protected simplification
gives 6,500 triangles, 51.85\% fewer than the previous revision and 85.27\% fewer than the defective
44,122-triangle engineering input. File size decreases from 675,084 to 325,084 bytes. All 3,148 groove,
flat-flange, and corner faces are protected, including all 512 flange-plane facets. Further trials at 6,300
and 6,200 triangles exceed the shape/quality guard and are not used. The result is one closed manifold
component of genus one with no detected self-intersections or zero-area facets, and unchanged bounds.

The previous closed volume is 1938412.694702 mm$^3$ and the final volume is 1946864.540816 mm$^3$. Filling
the recesses and applying the new simplification changes tungsten volume by +0.436019\%. The smooth,
unreduced reference has volume 1946881.004462 mm$^3$; simplification alone changes it by -0.00084564\%.
Bidirectional sampling against that reference gives maximum displacement 0.047853 mm, below the 0.05 mm
guard. These sampled limits apply after the authorized recess closure, rather than across the filled
cavities. The original engineering mass remains undefined.

Independent surface checks confirm 256 coplanar triangles at each flange face, with areas matching the
complete 128-sided annuli at local z = -625 and -605 mm. The inner face radii are 45 and 50 mm respectively,
adjoining the retained fillets; outer radius is 100 mm. Every outer flange surface facet belongs to one flat
face or the cylindrical rim, with zero recess-wall facets. All 1,134 independent containment probes pass,
including a dense grid just inside both flange faces, all former bolt and recess centres, tube groove
floors, the filled tube hole, and the central beam passage. The final STL and smooth reference reproduce
byte for byte from the archived raw input.

\begin{table}[htbp]
\centering
\caption{FT inner-shield smooth-flange revision.}
\label{tab:ft-shield-smooth-mesh}
\begin{tabular}{lrr}
\toprule
Quantity & Previous & Smooth \\
\midrule
Triangles & 13,500 & 6,500 \\
Bytes & 675,084 & 325,084 \\
Flange recesses & 26 & 0 \\
Boundary edges & 0 & 0 \\
Volume (mm$^3$) & 1938412.695 & 1946864.541 \\
\bottomrule
\end{tabular}
\end{table}

Two separate tracking comparisons use the previous closed 13,500-triangle mesh and the new smooth, unreduced
reference as baselines respectively. Both use the same fixed GEMC3 executable snapshot, beamline\_W
material, original placement, and matched 5 GeV target-origin primaries over 0.7--2.7 degrees and full
azimuth. FTFP\_BERT runs without optics or magnetic field, with one thread, default production cuts, and no
GUI or event output. Five alternating-order repetitions use seeds 12345--12349 and 200 events per electron,
proton, positive-pion, and neutron run. Each comparison has 52 processes including startup and isolated
overlap checks; all 104 complete normally and all eighty transport runs report 200 events, with no Geant4
exceptions in any fixture.

Previous/smooth median transport ratios are 1.273, 1.152, 1.160, 1.134 for electrons, protons, positive
pions, and neutrons respectively. These incremental upgrade timings include the added tungsten from filling
the blind recesses. Reference/smooth medians are 1.055, 1.023, 1.041, 1.030 in the same particle order,
comparing both cases with fully smooth flange surfaces and the recorded small decimation volume difference.
A ratio above one denotes faster tracking. Both comparisons are retained separately; neither establishes
full CLAS12 speedup or detector-response equivalence. Median startup simulation time changes from 0.776447
to 0.739538 s, with peak startup RSS 135.95 to 124.98 MiB.

\begin{table}[htbp]
\centering
\caption{Median transport wall seconds per 200 events: incremental upgrade.}
\label{tab:ft-shield-smooth-upgrade}
\begin{tabular}{lrrr}
\toprule
Particle & Previous & Smooth & Previous/smooth \\
\midrule
e- & 23.96760 & 18.82900 & 1.273 \\
proton & 2.23387 & 1.93947 & 1.152 \\
pi+ & 2.97310 & 2.56412 & 1.160 \\
neutron & 2.18361 & 1.92635 & 1.134 \\
\bottomrule
\end{tabular}
\end{table}

\begin{table}[htbp]
\centering
\caption{Median transport wall seconds per 200 events: smooth-reference control.}
\label{tab:ft-shield-smooth-reference}
\begin{tabular}{lrrr}
\toprule
Particle & Reference & Smooth & Reference/smooth \\
\midrule
e- & 20.41750 & 19.34670 & 1.055 \\
proton & 1.98474 & 1.94027 & 1.023 \\
pi+ & 2.68023 & 2.57448 & 1.041 \\
neutron & 1.96960 & 1.91241 & 1.030 \\
\bottomrule
\end{tabular}
\end{table}

The unchanged main-cone is included in a separate two-volume 10,000-point assembly check. Both previous and
smooth cases retain five pre-existing shield/cone overlap reports and one solid warning from the cone. Full
default construction completes without missing meshes, with one solid warning and eight overlap reports in
other or pre-existing joints. Smooth flange closure does not resolve that separate cone interference. All
fifteen native geometry/material and fourteen CAD variations continue to pass semantic comparison against
GEMC2. The earlier study archive is retained unchanged; this revision has a separate archive with its
previous mesh, smooth reference, final mesh, complete raw benchmarks, and validation evidence.

\FloatBarrier
\begin{figure}[htbp]
\centering
\includegraphics[width=0.9\textwidth]{figures/ft_shield_smooth_comparison.pdf}
\caption{Previous flange with blind recesses (left) and final smooth flange (right), with matched tube
sections showing both longitudinal grooves and the central bore retained.}
\label{fig:ft-shield-smooth}
\end{figure}
